(1) 对半径 \(r\) 的圆柱侧面，面积为 \(2\pi r\ell\)，稳态能量守恒要求同一 \(\dot Q=-2\pi r\ell\,k(T)dT/dr\) 对所有 \(r\) 恒定。因内热外冷，\(dT/dr<0\)，故 \(\dot Q>0\) 与正向约定一致。分离变量得到 \(k(T)dT=-(\dot Q/2\pi\ell)\,dr/r\)。这里的几何因子 \(r\) 会产生 \(\ln(b/a)\)，不能照搬平板公式中的厚度 \(b-a\)。

(2) 对 \(n\ne-1\)，积分 \(a\to r\) 有 \(k_0[T(r)^{n+1}-T_h^{n+1}]/[(n+1)T_0^n]=-\dot Q\ln(r/a)/(2\pi\ell)\)。在 \(r=b\) 使用 \(T_c\) 得 \(\dot Q=\{2\pi\ell k_0/[(n+1)T_0^n\ln(b/a)]\}(T_h^{n+1}-T_c^{n+1})\)。消去热流率后 \(T(r)^{n+1}=T_h^{n+1}-(T_h^{n+1}-T_c^{n+1})\ln(r/a)/\ln(b/a)\)。即使 \(n<-1\)，积分结果符号仍使热流为正，须整体保留分子与 \(n+1\) 的配对。

(3) 当 \(n=-1\)，导热率 \(k=k_0T_0/T\)，温度积分变为对数，结果 \(\dot Q=2\pi\ell k_0T_0\ln(T_h/T_c)/\ln(b/a)\)，且 \(\ln[T(r)/T_h]=-\ln(T_h/T_c)\ln(r/a)/\ln(b/a)\)。将一般公式的 \(T_h^{n+1}-T_c^{n+1}\) 除以 \(n+1\) 并令 \(n+1\to0\) 也得同一极限。对 \(n=0\) 得 \(\dot Q=2\pi\ell k_0(T_h-T_c)/\ln(b/a)\)，恰是常导热率同轴圆筒的热阻公式。

(4) 局域熵产生密度 \(\sigma=\mathbf j_q\cdot\nabla(1/T)=k(T)|\nabla T|^2/T^2\ge0\)，其中热流密度 \(\mathbf j_q=-k\nabla T\)。积分整个圆筒体积，用 \(\dot Q\) 与 \(dT/dr\) 消去温度梯度，得到 \(\dot S_{\rm gen}=\dot Q\int_a^b[-(1/T^2)dT/dr]dr=\dot Q(1/T_c-1/T_h)>0\)。边界热库熵率相加恰为该数：内热库流出热量给 \(-\dot Q/T_h\)，外冷库收到热量给 \(+\dot Q/T_c\)，与局域第二定律完全一致。
